\documentclass[12pt]{article}

% Core math packages
\usepackage{amsmath}
\usepackage{amssymb}
\usepackage{amsthm}
\usepackage{geometry}
\usepackage[utf8]{inputenc}
\usepackage{hyperref}
\usepackage{graphicx}
\usepackage[shortlabels]{enumitem}
\usepackage{tikz}
\usepackage{pgfplots}
\usepackage[most,skins]{tcolorbox}

% Custom styles shared with other HSC projects
\usepackage{styles/dl101-colors}
\usepackage{styles/dl101-hints}
\usepackage{styles/dl101-hsc-problems}

% TikZ libraries for diagrams
\usetikzlibrary{arrows.meta, calc, patterns, decorations.pathreplacing}
\pgfplotsset{compat=1.18}

\geometry{
  a4paper,
  total={170mm,257mm},
  left=20mm,
  top=20mm,
}

\author{Vu Hung Nguyen}
\title{HSC Math Extension 1 \& 2: Functions}
\date{}

\begin{document}

\maketitle

\begingroup\small\noindent This work is licensed under CC BY 4.0, see the LICENSE file on the github page for more info.\par\endgroup

\tableofcontents
\newpage

\section{Introduction}

\subsection{Project Overview}
This booklet compiles high-quality function problems curated specifically for the HSC Mathematics Extension 1 and Extension 2 syllabuses. 

\newpage
\section{Functions and Their Properties}
% Functions and their properties problems


\newpage
\section{The Inverse of a Function}
% The inverse of a function problems


\newpage
\section{Graphing the Inverse of a Function}
% Graphing the inverse of a function problems


\newpage
\section{Composite Functions}
\input{problems/04-Composite-functions}

\newpage
\section{Inverse Functions}
\input{problems/05-Inverse-functions}

\newpage
\section{Graphs of Inverse Functions}
% Graphs of inverse functions problems


\newpage
\section{Trigonometric Functions}
\input{problems/07-Trigonometric-functions}

\newpage
\section{Conclusion}

\subsection{Final Thoughts}
Mastering functions in HSC Mathematics requires practice and a deep understanding of their graphical properties, inverses, and algebraic manipulations.

\subsection{Contact Information}
\begin{itemize}[leftmargin=*]
  \item \textbf{LinkedIn:} \url{https://www.linkedin.com/in/vuhung16au/}
  \item \textbf{GitHub:} \url{https://github.com/vuhung16au}
  \item \textbf{Repository:} \url{https://github.com/vuhung16au/math-olympiad-ml/tree/main/HSC-Functions}
\end{itemize}

\end{document}
